\begin{table}[htbp]\centering\footnotesize
\caption{La auditoría de etiquetado de los anexos transversales}\label{cua:auditoria}
\begin{tabular}{>{\raggedright\arraybackslash}p{5.0cm}r>{\raggedleft\arraybackslash}p{2.2cm}>{\raggedleft\arraybackslash}p{2.8cm}}
\toprule
Indicador & 2026 & 2027, mismos anexos & 2027, con el 32 y el 33 \\
\midrule
Programas etiquetados & 226 & 228 & 243 \\
Programas en más de un anexo & 123 & 123 & 136 \\
\textbf{\% del etiquetado en programas con más de una etiqueta} & \textbf{95,4} & \textbf{95,8} & \textbf{95,2} \\
Suma aritmética de etiquetas (mdp) & 4.284.445,1 & 4.381.269,1 & 4.599.683,9 \\
\textbf{Pensiones no contributivas en el Anexo 13 (\%)} & \textbf{48,1} & \textbf{57,2} & \textbf{57,2} \\
\midrule Anexo 32: programas ya etiquetados en otro anexo & --- & --- & 16 de 23 \\
Anexo 32: mdp ya etiquetados en otro anexo & --- & --- & 130.912,2 de 175.045,5 \\
\midrule Anexo 33: programas ya etiquetados en otro anexo & --- & --- & 54 de 62 \\
Anexo 33: mdp ya etiquetados en otro anexo & --- & --- & 43.322,2 de 43.369,4 \\
\bottomrule
\end{tabular}
\\[2pt]\begin{minipage}{0.92\textwidth}\scriptsize Fuente: tablas de los anexos 13 a 19, 31, 32 y 33 de los proyectos de decreto, extraídas programa por programa y validadas contra el «Total general» de cada anexo (cobertura del 100~\% salvo el Anexo 13 de 2027, 99,99~\%, y el 31, 99,78~\%; el 32 no imprime total). \texttt{datos/ruta\_a/anexos\_auditoria.csv}. Pensiones no contributivas: claves S176, S286 y S316. \textbf{La suma de etiquetas no mide gasto.}\end{minipage}
\end{table}
